\cotefile{79}
\scanpage{205}
\mspage{105}

On l'obtient en relation avec la structure de graphe \inserted{$(S,A_e)$} sur $S$, associée au choix d'une base \inserted{de $\det M=\Lambda^2M$} ou ``jauge'' de $M$

\[
e\in\det(M)^* \tag*{(34)}
\]

par la \deleted{condition} formule

\[
A_e=\bigl\{\{s,t\}\in\mathscr P_2(S)\mid \inserted{\exists x,y\in M^*,\ \text{t.q.}}\ x\wedge y=e\ \deleted{\unsure{\dot x=s,\dot y=t}}\ \inserted{\dot x=s,\dot y=t}\bigr\} \tag*{(35)}
\]

(ne dépend que de la structure $\frac12$ symplectique de $M$ $\alpha=\{e,-e\}$ définie par $e$). Posant

\[
\left\{\begin{aligned}
&x\wedge y=\psi_e(x,y)\,e\\
&\psi_e\ \text{forme bilinéaire alternée non dég. définie par }e
\end{aligned}\right. \tag*{(36)}
\]

On voit donc que l'on a, pour $s\in S$, $s=\dot x$

\[
\operatorname{Et}_{A_e}(s)=
\{\deleted{\unsure{t\in S}}\inserted{\dot y}\mid\inserted{y\in M^*,}\ \psi_e(x,y)=\unsure{1|{\pm}1}\ \deleted{\unsure{\dot y=t}}\}. \tag*{(37)}
\]

Considérons alors l'application linéaire

\[
\left|\begin{aligned}
&\Phi_e:M\xrightarrow{\sim}M^\vee \qquad\text{caractérisée par}\\
&\langle x,\Phi_e(y)\rangle=\psi_e(x,y)\qquad\text{i.e.}\quad x\wedge y=\langle x,\Phi_e(y)\rangle\cdot e
\end{aligned}\right. \tag*{(38)}
\]

elle induit, par passage au quotient

\[
\deleted{\text{on a donc}}\quad \varphi_e:S\longrightarrow S^\vee,\qquad\text{et on aura, pour }s\in S \tag*{(39)}
\]

\[
\operatorname{Et}_{A_e}(s)=\{t\in S\mid\langle s,\varphi_e(t)\rangle\in R_{\unsure{1|A}}\},\quad\text{où} \tag*{(40)}
\]

\[
\left\{\begin{aligned}
&R_1\subset S\times S^\vee,\\
&\mathord{''}\ \{\deleted{\unsure{(s,t)}}\inserted{(\dot x,\dot x^\vee)}\}\in S\times S^\vee\mid \langle x,x^\vee\rangle=1\deleted{\unsure{t}}\}
=\{(s,s^\vee)\in S\times S^\vee\mid\inserted{\langle s,s^\vee\rangle=\dot 1}\}
\end{aligned}\right. \tag*{(41)}
\]

% slanted
\margin{où on introduit l'accouplement $(42)\quad S\times S^\vee\longrightarrow\Lambda/{\pm\mathrm{id}_\Lambda}$ par passage au quotient à partir de $M^*\times M^{\vee *}\longrightarrow\Lambda$.}

\deleted{On voit donc que les} \inserted{\unsure{Supposons|Supposant}} les structures de graphe $(S,A_e)$ \unsure{stricts|strictes}, i.e. qu'on a

\[
\left\{\begin{aligned}
&\forall(x,x')\in M^*\times M^*,\quad\text{la relation}\\
&\left[\begin{aligned}
&\forall y\in M^*\quad x\wedge y\in(\det M)^*\} \iff x'\wedge y\in(\det M)^*\\
&\hspace{38mm}\implies x'\wedge y\in\{\pm x\wedge y\}
\end{aligned}\right]\\
&\text{implique}\quad x'\in\pm x,\quad\deleted{\text{\unsure{i.e. que }} a\in\unsure{\mathrm{Fix}(M^*)},}\\
&\deleted{\exists\,\unsure{x,x'\in a}\ \text{et que }\unsure{\in}}
\end{aligned}\right. \tag*{(\unsure{43})}
\]

ce qui, pour $\Lambda$ $\frac12$ local, signifie simplement \unsure{qu'on|qu'un|qu'en}

% archive sheet number: 204
\scanpage{206}
n'est pas dans le cas octaédral généralisé --

\[
2\Lambda=\{0,2\},\quad 2\,1_\Lambda\ne0 \tag*{(44)}
\]

et dans le cas $\Lambda=\mathbb Z/n\mathbb Z$ signifie simplement que $n\ne4$), Alors il y a corr. $1$-$1$ entre étoiles de $(S,A_e)$, et éléments de $S^\vee$, via la relation $R$ de (41), et la structure de graphe $(S,A_e)$ est associée à la structure

\[
(S,S^\vee,R\subset S\times S^\vee,\varphi_e:S\xrightarrow{\sim}S^\vee). \tag*{(45)}
\]

Notons, maintenant que le groupe

\[
Z'_\Lambda=\Lambda^*/{\pm1} \tag*{(46)}
\]

opère sur $S$ et sur $S^\vee$ par symétrie, \deleted{et} par des opérations que je vais noter \deleted{$\unsure{T_\lambda,T^\vee_\lambda}$} \inserted{multiplicativement} ($\dot\lambda\in Z'_\Lambda$ \inserted{provenant de $\lambda\in\Lambda^*$})

\[
\deleted{\unsure{T_\lambda}}\dot\lambda\cdot\dot x\overset{\text{déf}}=(\lambda x)^{\boldsymbol\cdot},\qquad
\deleted{\unsure{T^\vee_\lambda}}\dot\lambda\cdot(x^\vee)^{\boldsymbol\cdot}\overset{\text{déf}}=(\lambda x^\vee)^{\boldsymbol\cdot}. \tag*{(47)}
\]

Ces \deleted{\unsure{``|«} \illegible{1} \unsure{''|»}} opérations sont des opérations \uline{libres}, et on a

\[
\left\{\begin{aligned}
&\forall(s,\unsure{t^\vee|s^\vee})\in S\times S^\vee,\ \text{\unsure{et|et à} }\lambda\in Z'_\Lambda,\ \text{on a}\\
&\boxed{(s,s^\vee)\in R\iff(\deleted{\unsure{T_\lambda}}\lambda s,\deleted{\unsure{T_{\lambda^{-1}}}}\lambda^{-1}s^\vee)\in R}.
\end{aligned}\right. \tag*{(48)}
\]

Notons aussi que

\[
\varphi_e=\varphi_{-e}:S\mathrel{\overset{\sim}{\rule{1.5em}{0.35pt}}}S^\vee, \tag*{(49)}
\]

\deleted{donc pour $\lambda\in\Lambda^*$, $\varphi_{\lambda e}$ ne dépend que de} i.e. $\varphi_e$ ne dépend que de l'élément

\[
\alpha=\{e,-e\}\in\Pi\overset{\text{déf}}=(\det M)^*/{\pm1}
\underset{e\mapsto e^{-1}}{\overset{\sim}{\rule{1.5em}{0.35pt}}}(\det M^\vee)^*/{\pm1} \tag*{(50)}
\]

% slanted
\margin{cf (54) provisoire}

défini par $e$, de sorte qu'on a en fait

\[
(\varphi_\alpha)_{\alpha\in\Pi},\quad
\begin{aligned}&\text{famille de bijections indexée}\\ &\text{par }\Pi\end{aligned} \tag*{(51)}
\]

Notons que \deleted{$\Pi$ est un torseur}

% slanted
\margin{cf (55)}

\[
\left\{\begin{aligned}
&\Pi\ \text{est un torseur sous }Z'_\Lambda\ \deleted{\unsure{\lambda}},\ \text{par des}\\
&\text{opérations notées }\alpha\mapsto T_{\dot\lambda}\alpha\ (\dot\lambda\in Z'_\Lambda),\ \text{déduites}\\
&\text{par passage au quotient de l'op. }e\mapsto\lambda.e.
\end{aligned}\right. \tag*{\text{(52) provisoire}}
\]

% archive sheet number: 205
\scanpage{207}
\mspage{106}

\deleted{\unsure{Je}} \deleted{On considère alors $\varphi_{\dot\lambda\dot e}$}

Notons que $\varphi_{\dot\lambda\dot e}(\dot x)\in S^\vee=M^{\vee *}/{\pm1}$ est défini \deleted{\unsure{via}} \inserted{comme} $\Phi_{\lambda e}(x)^{\boldsymbol\cdot}\deleted{\unsure{e}}$, où \deleted{$\varphi_{\lambda e}(x)$ défini par} \inserted{$\Phi_{\lambda e}(x)\in M^\vee$ défini par (38)}

\[
(y,\Phi_{\lambda e}(x))\cdot(\lambda e)=y\wedge x
\]

qu'on compare avec

\[
(y,\Phi_e(x))\,e=y\wedge x
\]

en comparant on trouve

\[
\Phi_{\lambda e}(x)=\lambda^{-1}\Phi_e(x)=\Phi_e(\lambda^{-1}x)
\]

soit, en passant aux éléments correspondants de $S^\vee$

\[
\varphi_{\dot\lambda.\alpha}(s)\deleted{\unsure{\varphi_{T_\lambda\alpha}(x)}}=
\dot\lambda^{-1}\varphi_e(s)=\varphi_e(\dot\lambda^{-1}s)
\qquad\forall s\in S^\vee,\ \dot\lambda\in Z'_\Lambda. \tag*{(53)}
\]

\margin{\uline{provisoire} cf (55)}

Si on avait utilisé $\Pi^\vee=(\det M^\vee)^*$ pour indexer les $\varphi:S\longrightarrow S^\vee$ que nous \unsure{introduisons|induisons|considérons}, ce qui était plutôt plus naturel (puisque l'application lin. \deleted{$M\longrightarrow M^\vee$ qui} $\Phi:M\mathrel{\rule{1.6em}{0.35pt}}M^\vee$ définie au signe près, qui relève $\varphi$, peut \deleted{\unsure{se}} directement être considérée comme une application bilinéaire (alternée) $M\times M\longrightarrow\Lambda$, i.e. une forme linéaire $\Lambda^2M\mathrel{\rule{1.6em}{0.35pt}}\Lambda$, et s'identifie donc, ce titre à un élément $\alpha^\vee$ ($=\alpha^{-1}$) de.

\[
(\det M^\vee)^{\prime}\simeq(\det M)^\vee
\]

Aussi on va considérer plutôt qu'on obtient

\[
\Pi_{\deleted{\unsure{M}}}=\Pi_M=(\det M^\vee)^*/{\pm\mathrm{id}_{M^\vee}} \tag*{(54)}
\]

sur lequel on fait opérer \deleted{$\unsure{\Lambda}$} $Z'_\Lambda=\Lambda^*/{\pm1}$ par l'opération habituelle -- qui \deleted{\unsure{s'oppose}} \inserted{correspond} à la multiplication d'une forme sur $M$ par un scalaire -- et (53) \deleted{\unsure{sera}} remplacée par

\[
\boxed{\varphi_{\lambda.\alpha}(s)=\dot\lambda\varphi_\alpha(s)=\varphi_{\unsure{\alpha|s}}(\dot\lambda s)}
\quad\begin{aligned}
&\forall\alpha\in\Pi=(\det M^\vee)^*/{\pm1}\\[-1pt]
&\lambda\in Z'_\Lambda=\Lambda^*/\{\pm1\}
\end{aligned} \tag*{(55)}
\]

% archive sheet number: 206
\scanpage{208}
\deleted{\unsure{Le|4.}} Appelons ``prégraphe'' \unsure{un|une} système

\[
\Gamma=(S,S^\vee,R\subset S\times S^\vee) \tag*{(56)}
\]

de deux ensembles $S,S^\vee$, et d'une relation entre eux -- il sera dit \uline{séparant} si les applications

\[
S\mathrel{\unsure{\rule{1.35em}{0.35pt}|\longrightarrow}}\mathscr P(S^\vee),\qquad S^\vee\mathrel{\unsure{\rule{2em}{0.35pt}|\longrightarrow}}\mathscr P(S)
\]

correspondantes sont injectives. Le symétrique du prégraphe (56) est le prégraphe

\[
\Gamma^\vee=(S^\vee,S,{}^sR\subset S^\vee\times S). \tag*{(57)}
\]

On appelle ``\uline{homothétie}'' du prégraphe toute paire de permutations

\[
\lambda=(\lambda_S,\lambda_{S^\vee})\in\unsure{\mathfrak S|\mathscr S}_S\times\mathfrak S_{S^\vee} \tag*{(58)}
\]

satisfaisant la condition

\[
\left\{\begin{aligned}
&(\lambda_S,\lambda_{S^\vee}^{-1})(R)=R\qquad\text{i.e.}\\
&\forall(s,s^\vee)\in S\times S^\vee,\ \text{on a }\bigl((s,s^\vee)\in R\iff(\lambda s,\lambda^{-1}s^\vee)\in R,
\end{aligned}\right. \tag*{(59)}
\]

\deleted{\unsure{Ce|Cela|elles}} sont donc les éléments \deleted{de $(\lambda,\lambda')$} \deleted{$\unsure{(s,t)\mapsto}$} $(u,u')\in\mathfrak S_S\times\mathfrak S_{S^\vee}$ tels que l'on ait

\[
\left\{\begin{aligned}
&(u,u'^{-1})\in\operatorname{Aut}(S,S',R)\qquad\text{i.e.}\\
&(u^{-1},u')\in\operatorname{Aut}(S,S',R)
\end{aligned}\right. \tag*{(60)}
\]

Notons que si $R$ est séparante en $S^\vee$, alors $\lambda_{S^\vee}$ est \unsure{connu|connue} quand on connaît $\lambda_S$, \unsure{celui|qui|elle|et} est défini à partir de $\lambda_S$ par

\[
\deleted{\unsure{\lambda_S(s^\vee)|\lambda_{S^\vee}(s^\vee)}}\
\left\{\begin{aligned}
&{}^sR(\lambda_{S^\vee}(s^\vee))=\lambda_S^{-1}(R(s^\vee))\qquad\text{ou aussi}\\
&{}^sR(\lambda_{S^\vee}^{-1}(s^\vee))=\lambda_S(R(s^\vee))
\end{aligned}\right. \tag*{(61)}
\]

et inversement, si $R$ est sép en $S$, $\lambda_{S^\vee}$ détermine $\lambda_S$.

% archive sheet number: 207
\scanpage{209}
\mspage{107}

donc dans le cas biséparant, $\lambda_S$ et $\lambda_{S^\vee}$ se déterminent mutuellement. Dans ce cas, on trouve des bijections canoniques

\[
\left\{\begin{aligned}
&\begin{gathered}\text{Homothéties de }\deleted{\unsure{R}}\Gamma\\
\text{i.e. couples }\lambda_S,\lambda_{S^\vee}\\\text{satisfaisant (59)}\end{gathered}
\simeq\overbrace{\begin{gathered}\text{automorphismes}\\\text{de la structure}\\
\Gamma=\underset{s}{(}S,S^\vee,R)\end{gathered}}^{\text{soit }Z}\\[3pt]
&\underset{\text{soit }Z_S}{\simeq}
\bigl\{\lambda\in\unsure{\mathfrak S|\mathscr S}_S\mid
\underbrace{\lambda(\mathcal E)}_{\substack{\text{transport}\\\text{de str}}}=\mathcal E\bigr\}
\quad\begin{gathered}\text{où }\mathcal E\subset\mathscr P(S)\\ \text{est l'image de}\\ S^\vee\xrightarrow{{}^sR}\mathscr P(S)\end{gathered}\\
&\underset{\text{soit }Z_{S^\vee}}{\simeq}
\bigl\{\lambda^\vee\in\mathfrak S_{S^\vee}\mid
\underbrace{\lambda^\vee(\mathcal E^\vee)}_{\substack{\text{transp.}\\\text{de str.}}}=\mathcal E^\vee\bigr\}
\quad\begin{gathered}\text{où }\mathcal E^\vee\subset\mathscr P(S)\\ \text{est l'image de}\\S\xrightarrow R\mathscr P(S^\vee)\end{gathered}
\end{aligned}\right. \tag*{(62)}
\]

Sauf le premier des trois \inserted{4} ensembles, les autres ensembles sont munis des structures de groupes évidentes, \deleted{qui} sont isomorphes par les applications évidentes

\diagram{\((63)\quad\)\begin{tikzcd}[ampersand replacement=\&,row sep=1.5em,column sep=1.5em]
\& Z\arrow[dl,"{\substack{p_S\\[-2pt]\sim}}"']\arrow[dr,"{\substack{p_{S^\vee}\\[-2pt]\sim}}"] \&\\
Z_S\arrow[d,hook] \& \& Z_{S^\vee}\arrow[d,hook]\\
\mathfrak S_S \& \& \mathfrak S_{S^\vee}
\end{tikzcd}}

mais l'ens des homothéties apparaît comme le sous-ensemble un peu vicieux de $\mathfrak S_S\times\mathfrak S_{S^\vee}$

\[
\deleted{\unsure{Z}}\ \mathcal H=\{(p_S(\lambda),p_{S^\vee}(\lambda^{-1}))\in\mathfrak S_S\times\mathfrak S_{S^\vee}\mid\lambda\in Z\}, \tag*{(64)}
\]

qui fait que ce \uline{n'est pas} \inserted{a priori} un sous-groupe.

% archive sheet number: 208
\scanpage{210}
Si p. ex. $R$ est le graphe d'une \uline{bijection} entre $S$ et $S^\vee$, alors \deleted{\unsure{on|on a}} $Z_S=\unsure{\mathfrak S|\mathscr S}_S$, $Z_{S^\vee}=\mathfrak S_{S^\vee}$, et identifiant $S^\vee$ à $S$ par $R$, \inserted{d'où} $\mathcal H\subset\mathfrak S_S\times\mathfrak S_S$, $\mathcal H$ s'identifie à l'antidiagonale

\[
\mathcal H=\{(u,u^{-1})\in\mathfrak S_S\times\mathfrak S_S\mid u\in\mathfrak S_S\}, \tag*{(}
\]

qui n'est pas un \unsure{ss-groupe|ss-gpe|sous-groupe} \deleted{si $\operatorname{card}S\ge3$} \unsure{ssi.|dès que} $\mathfrak S_S$ n'est pas comm i.e. $\operatorname{card}S\ge3$. La définition \inserted{de $\mathcal H$} peut paraître saugrenue, pourtant elle est fortement suggérée par les considérations précédentes.

Ce n'est pas simplement un canular de notation -- ainsi, rien ne nous empêcherait de modifier la définition en apparence saugrenue des couples $(\lambda,\lambda^\vee)$ qualifiés d'homothéties, en les remplaçant par les couples $(\lambda^{-1},\lambda^\vee)$ ou $(\lambda,\lambda^{\vee-1})$ -- par les automorphismes du graphe, ce qui reviendrait à changer de notation pour les relations entre $\varphi_0,\varphi$ via $\lambda,\lambda^\vee$, en écrivant

\[
\deleted{\unsure{\lambda}}\
\left\{\begin{aligned}
&\varphi=\lambda^\vee\varphi_0=\varphi_0\lambda^{\unsure{\vee-1|-1}}
[=(\lambda\varphi_0^{-1})^{-1}]\\
&\deleted{\text{donc}}\quad\deleted{\varphi_0(\lambda}
\left\{\begin{aligned}&\lambda^\vee=\varphi_0(\lambda)^{-1}\\&\lambda=(\varphi_0^{-1}(\lambda^\vee))^{-1}\end{aligned}\right.\\
&\varphi_0(\lambda^{-1}s)=\lambda^\vee\varphi_0(s)\quad\text{ou }\varphi_0(\lambda s)=\lambda^{\vee-1}\varphi_0(s)
\end{aligned}\right.
\]

% slanted
\margin{Cette définition de $\lambda,\lambda^\vee$ en termes de $\varphi_0,\varphi$ a aussi l'avantage d'être stable par symétrie, qui échange alors $\lambda,\lambda^\vee$.}

ce qui donne par contre

\[
(\lambda,\lambda^\vee)(R)=R\quad\text{i.e.}\quad
\inserted{\forall s,s^\vee\in S\times S^\vee}\ (s,s^\vee)\in R\iff(\lambda s,\lambda^\vee s^\vee)\in R.
\]

Dans le cas de l'\unsure{exemple|example} linéaire, ça signifie que si on fait opérer les scalaires \inserted{$\in\Lambda^*$} sur $M$ de la façon habituelle, on les fait opérer sur $M^{\unsure{\vee|*}}$ par contragrédience (la ``multiplication'' inhabituelle), ce qui a l'avantage

% archive sheet number: 209
\scanpage{211}
\mspage{108}

de respecter le produit scalaire \unsure{donc|dans} la relation d'incidence $R$, mais l'inconvénient que les applications ne commutent plus à l'op de $\Lambda^{\unsure{\alpha|\times|*}}$ sur \deleted{$\unsure{M}$} $S$ et $S^\vee$, mais qu'il \deleted{faut} écrire $\varphi_0(\lambda s)=\lambda^{\unsure{\vee-1|-1}}\varphi_0(s)$ (au lieu de $\lambda\varphi_0(s)$). L'avantage d'avoir une loi de groupe bien en évidence sur les homothéties, dans cette notation, peut paraître prépondérant. Mais il faut bien noter qu'en général, \uuline{on n'a pas} d'opération du groupe $Z$ sur l'ens $\Pi$ des ``jauges'' du prégraphe $\Gamma$.

\deleted{On a bien défini une application}

\[
\begin{aligned}
&\deleted{Z_{\unsure{s|S}}\times\Pi\longrightarrow\Pi}\\
&\deleted{(\lambda,\varphi_0)\longmapsto\varphi_0\lambda^{-1}=\lambda^\vee\varphi_0,}
\end{aligned}
\]

Pour que $(\lambda,\lambda^\vee)\in\unsure{\mathfrak S|\mathscr S}_S\times\mathfrak S_{S^\vee}$ définisse \deleted{une} -- à partir de \deleted{\unsure{à}} une jauge $\varphi_0\in\Pi$ -- une deuxième jauge $\varphi$ \inserted{$=\lambda^\vee\varphi_0=\varphi_0\lambda^{-1}$}, il faut en effet, en plus de la condition que $(\lambda,\lambda^\vee)$ soit un autom du prégraphe $\Gamma$, la condition (\uline{qui dépend du choix de $\varphi_0$})

\[
\lambda^\vee\varphi_0=\varphi_0\lambda^{-1}\quad\text{i.e.}\quad\varphi_0=\lambda^\vee\varphi_0\lambda. \tag*{(65)}
\]

Si on décide d'oublier $\lambda^\vee$ (p. ex.) pour ne retenir que $\lambda$ (et définir $\varphi$ par $\varphi_0\lambda^{-1}$, \inserted{et $\dot\lambda^\vee$ par $\lambda^\vee=\varphi_0^{-1}(\lambda^{-1})$ \deleted{\unsure{p.}}}), on est confronté à ce \deleted{\unsure{spectacle}}\unsure{problème|première}, que la condition sur $\lambda$ qui exprime le \deleted{symétrie} fait que

\[
(\lambda,\lambda^\vee)=(\lambda,\varphi_0^{-1}(\lambda^{-1}))
\]

\deleted{satisfasse à \unsure{une} la condition de symétrie, (} invarie $R$ \unsure{1)|1\textsuperscript{o}} dépend de $\varphi_0$, et 2\textsuperscript{o}) ne définit pas à priori un \uline{sous-groupe} de $\mathfrak S_S$ ! En fait, si on exprime

% archive sheet number: 210
\scanpage{212}
cette condition sous la forme (29'),

\[
\unsure{(*)}\quad\forall s,t\in S,\quad\{s,t\}\in A_0\iff\{\lambda^{-1}s,\lambda t\}\in A_0,
\]

on voit ceci

\[
\begin{aligned}
&\text{Si }\lambda,\mu\in\unsure{\mathfrak S|\mathscr S}_S\ \text{satisfont séparément à la condition}\\
&\unsure{(*)|\text{cas lin.}}\ \text{(i.e. sont des homothéties pour le graphe }(S,A_0)),\\
&\text{alors }\lambda\mu\ \text{y satisfait \unsure{sss|ssi|si} on a }\lambda\mu=\mu\lambda\\
&\deleted{\text{\unsure{sur}}}\quad\text{(cas d'un graphe \uline{séparant}).}
\end{aligned} \tag*{(66)}
\]

Pour continuer le travail, on supposera donc donné un sous-groupe \uline{commutatif}

\[
\deleted{\unsure{Z}}\qquad Z\subset\mathfrak S_S
\]

formé d'homothéties de $(S,A_0)$, et on considère \inserted{alors} le \deleted{éléments} \inserted{sous-groupe} correspondant \inserted{$Z_\Gamma$} de \deleted{$\unsure{\mathfrak S_S}\ \mathfrak S_S\times\mathfrak S_{S^\vee}$} $\operatorname{Aut}(S,S^\vee,R)$, on trouve : \deleted{que les}

\[
\left\{\begin{aligned}
&\text{les }\varphi_{\unsure{a|\alpha}}:S\longrightarrow S^\vee\ \text{telles qu'on ait}\\
&\lambda^\vee\varphi_a\lambda=\varphi_a\quad\text{pour tout }\lambda\in Z_\Gamma\subset\operatorname{Aut}_S(\Gamma)\\
&\text{forment un }\deleted{\text{torseur sous le groupe }Z_\Gamma,}\ \inserted{\unsure{Z_\Gamma|\Gamma}\text{-ens }\Pi_Z,}\ \text{par}\\
&\text{par l'opération }\deleted{\text{\unsure{qui à}}}\ \text{du groupe }Z_\Gamma,\ \text{qui}\\
&\hspace{48mm}\text{y opère \uline{librement} par}\\
&T_{\lambda,\lambda^\vee}\varphi=\lambda^\vee\varphi=\varphi\lambda^{-1}.
\end{aligned}\right. \tag*{(67)}
\]

% slanted
\margin{On a supposé que $\Pi_Z\ne\varnothing$, car on est parti d'une $\varphi_0\in\Pi_Z$.}

En effet, si \inserted{$(\lambda,\lambda^\vee)\in Z$,} on doit vérifier \inserted{d'abord} que

\[
\inserted{\varphi\in\Pi_Z}\Longrightarrow\underbrace{T_{\lambda,\lambda^\vee}\varphi}_{\lambda^\vee\varphi}\in Z\quad\text{i.e.}\quad
\forall(\mu,\mu^\vee)\in Z,\ \deleted{\text{qu'}}\text{on a}
\]

\[
\begin{aligned}
&\underbrace{\mu^\vee(\lambda^\vee}_{\lambda^\vee\mu^\vee}\varphi)\mu=(\lambda^\vee\varphi)\\
&\hspace{18mm}\underbrace{\phantom{\mu^\vee\varphi\mu}}_{\varphi}
\end{aligned}
\]

O,K, Le fait que $Z_\Gamma$ opère librement est clair, \deleted{\unsure{pour}} car pour $\varphi_0,\varphi$ \unsure{fixés|fixées}, il n'y a qu'un seul $\lambda,\lambda^\vee$ \ldots

\deleted{On \unsure{fait} le fait} \deleted{ensuite que si $\varphi_0\in$} On aimerait que

% archive sheet number: 211
\scanpage{213}
% archive sheet number: 212

\mspage{109}

cette opération \unsure{ci|de $Z$} soit transitive -- soit donc $\varphi_0\in\Pi_Z$, \unsure{déjà|disons}, et $\varphi\in\Pi_Z$, on peut écrire

\[
\varphi=g_0^\vee\varphi_0=\varphi_0g_0^{-1}\qquad\text{où }(g_0,g_0^\vee)\in\operatorname{Aut}(\Gamma),
\]

\inserted{\uline{est-il dans $Z_\Gamma$ ?}} uniquement déterminé, on \deleted{doit} exprimer la

\[
\text{condition}\qquad\varphi\in\Pi_Z\quad\text{i.e.}
\]

% slanted
\margin{\uline{Lemme} Soit $\varphi_0\in\Pi_Z$ une $Z$-jauge, et $g\in\operatorname{Aut}(\Gamma)$, telle que l'on ait \deleted{\unsure{$g$}} $\underbrace{g_S^\vee\varphi_0=\varphi_0g_S^{\unsure{-1|{}}}}_{\varphi}$ (i.e. $(g_S,g_S^\vee)$ définit une homothétie de la structure de graphe $A_{\varphi_0}$ associée à $\varphi_0$). Conditions équivalentes : a) $g_S^\vee\varphi'=\varphi'g_S^{-1}$ pour \uline{tt} $\varphi'$ de la forme $T_\lambda\varphi_0$ ($\lambda\in Z$) i.e. $g_S$ est une homothétie \inserted{pour \uline{ttes} les str. de graphes $A_\varphi$} \inserted{$\leftarrow$ \uline{tt} $\varphi'\in\Pi$} b) $\varphi$ ($\overset{\text{déf}}=g_S^\vee\varphi_0=\varphi_0g_S^{-1}$) est également une $Z$-jauge c) $g\in\operatorname{Cent}(Z)$}

\[
\underbrace{\lambda^\vee\varphi\lambda}=\underbrace{\varphi},\quad\forall(\lambda,\lambda^\vee)\in Z\quad\text{i.e., \unsure{écrivant|écrivons} }\varphi=\varphi_0g_0^{-1}
\]

\[
\deleted{\unsure{\lambda^\vee g_0^\vee\varphi_0\lambda=g_0^\vee\varphi_0}}\quad\unsure{2}
\]

\[
\begin{aligned}
&\deleted{\unsure{g_0}}\quad\lambda^\vee\varphi_0g_0^{-1}\lambda=\varphi_0g_0^{-1}\qquad\text{i.e.}\\
&\underbrace{\lambda^\vee\varphi_0}_{\varphi_0\lambda^{-1}}g_0^{-1}(\lambda)=\varphi_0\qquad\text{i.e.}\\
&\deleted{\unsure{R}}\qquad\deleted{\varphi_0=\lambda_0}\\
&\lambda^{-1}g_0^{-1}(\lambda)=1\\
&\lambda=g_0^{-1}(\lambda)\qquad\text{i.e. } g_0\ \deleted{\text{normalise}}\ \inserted{\text{centralise}}\ \text{tout }\lambda\in\unsure{\mathfrak S|\mathscr S}_S\\
&\text{ou }g_0(\lambda)=\lambda\qquad\text{provenant de }\deleted{\unsure{\mathfrak S_S}}\ Z
\end{aligned}
\]

\deleted{donc} on trouve

\[
\left\{\begin{aligned}
&\text{l'ens. }\Pi_Z\ \text{des jauges de }\Gamma=(S,S^\vee,R)\\
&\text{compatibles avec le groupe d'autom.}\\
&Z\subset\operatorname{Aut}(\Gamma),\ \text{i.e satisfaisant }\lambda^\vee\varphi\lambda\quad\forall(\lambda,\lambda^\vee)\in Z,\\
&\text{est un torseur sous le \uline{centralisateur}}\\
&\text{de }Z\ \text{dans }\operatorname{Aut}(\Gamma),\ \text{opérant par}\\
&T_g\varphi_0=g_{S^\vee}\varphi_0=\varphi_0g_S^{-1}
\end{aligned}\right. \tag*{(68)}
\]

\margin{\uline{faux}}

% slanted
\margin{Il faut écrire la petite vérification que si on a $g_S^\vee\varphi_{\unsure{a|0}}=\varphi g_S^{-1}$ pour \uline{un} $\varphi\in\Pi$, on l'a pour tous (\unsure{ceci pour|dès que} $g\in\unsure{\mathcal G|G}$).}

Ouf -- il aura fallu de la sueur, pour accoucher de cette petite souris !

\unsure{tout ce|tout-ce|tendre} qu'on a montré c'est que l'application

\[
\deleted{\unsure{\varphi}}\ g\longmapsto g_{S^\vee}\varphi_0=\varphi_0g_S^{-1}
\]

est une \uline{bijection} du \deleted{le} sous-ensemble $\mathcal G(\varphi_0)$ de $\mathcal G$, formé des $g$ satisfaisant $g_{S^\vee}\varphi_0=\varphi_0g_S^{-1}$ i.e. qui sont des $A_{\varphi_0}$-homothéties (\uline{et qui ne forment pas \inserted{néc.} un groupe} !) avec l'ens $\Pi_Z$ des $Z$-jauges !

% slanted
\margin{\uline{NB} on a $\mathcal G(\varphi_0)=\{g\in\operatorname{Aut}(\Gamma)\mid\varphi_0(g_S)g^\vee=1\}$ i.e. $g^\vee=\varphi_0(g_S^{-1})$}

% slanted
\margin{Il est donc prudent d'étendre un voile pudique sur l'ens $\Pi_Z$, qui n'est pas bien compris et peut-être \unsure{déshérité|d'intérêt|dénué d'intérêt}, pour se borner décidément, aux ens $\Pi$}

\uline{Mais il est vrai que $Z$ opère librement sur $\Pi_Z$}
\scanpage{214}
% slanted
\margin{\uline{Scholie}}

En résumé, nous \deleted{\unsure{sommes ainsi|sommes ici}} amenés à nous intéresser à la structure suivante

\[
\left\{\begin{aligned}
&\text{a) Triple }\Gamma=(S,S^{\unsure{*|\vee}},R)\ \text{de deux ens }S,S^\vee\ \text{et d'une}\\
&\qquad \text{relation }R\subset S\times S^\vee\\
&\text{b) Donnée d'un ss-groupe \uline{commutatif}}\\
&\hspace{20mm}Z\subset\operatorname{Aut}(\Gamma),\\
&\qquad\text{dont le centralisateur }\inserted{\text{dans }\operatorname{Aut}(\Gamma)}\ \text{est noté }\unsure{\mathcal G|G}.
\end{aligned}\right. \tag*{(69)}
\]

% slanted
\margin{dans le cas de l'exemple linéaire avec $M^*,M^{\vee *}\ldots$, $\mathcal G$ contient en tous cas le groupe $\operatorname{Gl}_\Lambda\bigl(\mathcal O(\unsure{*,M|M})\bigr)'$}

On appelle \uline{$Z$-jauges} d'une telle structure les \deleted{applications} bijections

\[
\left\{\begin{aligned}
&\varphi:S\xrightarrow{\sim} S^{\unsure{*|\vee}}\qquad\text{satisfaisant}\\
&1)\quad\deleted{\unsure{\varphi(s,t)}}(\mathrm{id}\times\varphi)^{-1}(R)\ \text{est une relation \uline{symétrique}}\\
&\qquad\text{sur }S\\
&\qquad\text{i.e. }\forall s,t\in S,\quad(s,\varphi(t))\in R\iff(t,\varphi(s))\in R\\
&2)\quad\forall(\lambda,\lambda^\vee)\in Z,\ \text{on a}\\
&\qquad\lambda^\vee\varphi\lambda=\varphi\qquad\text{i.e }\lambda^\vee\varphi=\varphi\lambda^{-1}
\end{aligned}\right. \tag*{(70)}
\]

Si $R$ est biséparante, de sorte que \deleted{$S^\vee$ soit} la donnée de $R$ s'identifie à celle de $(S,S^\vee\subset\mathscr P(S))$, alors les $\varphi$ satisfaisant (1) s'identifient aux structures de graphe \deleted{\unsure{sép}} \inserted{strict} séparant sur $S$ (sans exclure des boucles) $A_0\subset\mathscr P_2(S)$, qui admettent $S^\vee$ comme ens de leurs étoiles, et pour lesquels les $\lambda\in Z_S$ (provenant des $\lambda\in Z$) sont des \uline{homothéties}. \inserted{En tous cas,} Les $\varphi$ en question, \inserted{(s'il y en a !)}, forment un \uline{torseur} (non sous $Z$, mais) sous $\mathcal G=\operatorname{Cent}_{\operatorname{Aut}(\Gamma)}(Z)$.

Les $Z$-multigraphes envisagés précédemment sont donc les sous \inserted{$Z$-}torseurs de celui-ci, \deleted{i.e des} \inserted{donc s'identifient} aux points de $\mathcal G/Z$.

% slanted
\margin{canular, cf annotation page précédente. Mais $Z$ opère librement sur $\Pi_Z$}

% archive sheet number: 213
\scanpage{215}
\mspage{110}

Dans ce contexte général, si on part d'une ``vraie'' structure de graphe (sans boucles, i.e. $A_0\subset\mathscr P_2(S)$) séparé, il n'est pas clair qu'en \unsure{le|la} transformant par un élément de $Z$, faisant

\[
\text{donc}\quad A_\lambda=\{\{s,\lambda t\}\in\mathscr P(S)\mid s,t\in S\ \text{t.q.}\ \{s,\unsure{1|t}\}\in A_0\},
\]

on trouve encore des vrais graphes -- car il pourrait arriver qu'on ait $t=\lambda s$ pour un \deleted{$\lambda$} triple convenable $(s,t,\lambda)\in S\times S\times Z$, avec $\{s,t\}\in A_0$.

Prenons p. ex. cas $S^\vee=S$, $R={}$ diagonale, donc les jauges (satisfaisant 70 (1)) sont les $\varphi\in\unsure{\mathfrak S|\mathscr S}_S$ telles que la relation \deleted{$\unsure{(s,\varphi\in)}$} $\varphi(t)=s$ en $s,t$ soit symétrique, i.e. équivalente à $\varphi(s)=t$ i.e. à $s=\varphi^{-1}t$, ce qui signifie aussi

\[
\varphi^2=\mathrm{id}\qquad\text{i.e. }\varphi\text{ est une involution} \tag*{(71)}
\]

Les autom. \deleted{de} $(\lambda,\lambda^\vee)$ de $\Gamma=(S,S^\vee,\unsure{\mathrm{Diag}|R_{\mathrm{diag}}})$ sont les couples $(\lambda,\lambda)$ avec $\lambda\in\mathfrak S_S$, et la condition 70 (2) s'écrit

\[
\lambda\varphi\lambda=\varphi\qquad\text{soit}\quad\varphi(\lambda)=\lambda^{-1} \tag*{(72)}
\]

i.e. l'involution ``renverse'' $\lambda$. Si donc on se donne un sous-groupe commutatif $Z$ de $\mathfrak S_S$, alors $\Pi_Z$ est formé des involutions \inserted{$\varphi\in\mathfrak S_S$} qui \deleted{\unsure{opè}} normalisent $Z$, en y induisant l'autom. \deleted{$\unsure{Z}$} $\lambda\longmapsto\lambda^{-1}$. Si p. ex. $S$ est \inserted{l'ens. sous-jacent à} un torseur $L$ sous $Z$, \inserted{(donc $Z$ est son propre centralisateur dans $\mathfrak S_S$),} \deleted{ce sont les $\varphi$ en question s'identifient} \deleted{avec involutions} alors choisissant une origine dans $L$ d'où $L\simeq Z$, ce sont les involutions de la forme

\[
\varphi_\lambda(z)=\lambda-z\qquad \lambda\in Z,\ z\in Z,
\]

Ici la relation \deleted{$\unsure{z\in}\ s=\varphi_a(s)\ \text{i.e. }z=a-z$}

% archive sheet number: 214
\scanpage{216}
on a

\[
(z,\varphi_\lambda(z))\in R\iff\varphi_\lambda(z)=z\quad\text{i.e.}\quad2z=\lambda \tag*{(73)}
\]

et \deleted{rien n'empêche qu'il y ait} \inserted{il y a} des solutions \inserted{en $z$} de cette équation pour certaines valeurs du paramètre (\deleted{\unsure{notamment}} \inserted{savoir} pour $\lambda\in2Z$ -- NB on écrit additivement ici la loi de groupe), alors qu'il est possible qu'il n'y en ait pas pour d'autres (dès lors que $2Z\ne Z$) -- i.e. qu'on peut \inserted{néanmoins} choisir une ``origine'' $\varphi_{\lambda_0}$ telle que elle \unsure{si|se} corresponde à une ``vraie'' structure de graphe $A_{\lambda_0}$. L'exemple le plus économique serait \deleted{$\unsure{E}$} $Z=\mathbb Z/2\mathbb Z$ \unsure{,} \ldots

Soit $\Pi$ un sous-$Z$-torseur du $\mathcal G$-torseur des $Z$-jauges du prégraphe $\Gamma=(S,S^\vee,\Gamma)$. Ainsi $\unsure{\Pi|\Gamma}$ est formé de bijections

\[
\varphi:S\mathrel{\overset{\sim}{\rule{2em}{0.35pt}}} S^\vee
\]

satisfaisant

\[
\lambda_{S^\vee}^{\unsure{*|\vee}}\varphi(s)=\deleted{\lambda_S}\varphi(\lambda_S^{-1}(s))
\qquad \lambda_{S^\vee}^{*}\varphi=\varphi\circ\lambda_S^{-1} \tag*{(74)}
\]

i.e. compatibles avec les opérations de $Z$ sur $S$ et $S^\vee$, via l'hom.

\[
\begin{aligned}
&\lambda\longmapsto\lambda^{-1}\\
&Z\longrightarrow Z
\end{aligned} \tag*{(75)}
\]

des groupes d'opérateurs. \deleted{On peut donc considérer que $S^\vee$ est déduit de $S$ par le torseur} Si $\lambda\in Z$, on aura en tordant \deleted{\unsure{par}} des \unsure{rels|deux} \deleted{par le $Z$-torseur $\Pi$, pour l'opération de $Z$ sur $S$ $T_\lambda\varphi=\lambda_{S^\vee}\varphi=\varphi\lambda_S^{-1}$}

% archive sheet number: 215
\scanpage{217}
\mspage{111}

Considérons l'application

\[
\begin{aligned}
F:\ \Pi\times S&\longrightarrow S^\vee\\
(\varphi,s)&\longmapsto\varphi(s)
\end{aligned}
\]

on aura donc

\[
\underbrace{F(T_\lambda\varphi,s)}_{(\varphi\circ\lambda^{-1})(s)=\varphi(\lambda^{-1}s}
=F(\varphi,\lambda_S^{-1}s)
\]

ce qui montre que $F$ se factorise en une application canonique, qui est en fait iso d'ens.

\[
\Pi\wedge_Z(S!)\xrightarrow{\sim} S^\vee\qquad\varphi\wedge s\longmapsto\varphi(s) \tag*{(76)}
\]

où $S!$ désigne le $Z$-ensemble déduit de $S$ en y remplaçant l' \deleted{\unsure{structure|action}} \inserted{opération} donnée \inserted{``à gauche''} de $Z$ par l'opération ``à droite'' que s'en déduit par $\lambda\longmapsto\lambda^{-1}$.

\deleted{Comme les applications $\varphi$ commutent aux opérations de $Z$ sur $S$ et $S^\vee$} \inserted{Dans cet isomorphisme, on aura}

\[
\deleted{\unsure{et}}\quad
\underbrace{\lambda.(\varphi\wedge s)}_{\substack{(T_\lambda\varphi)\wedge s=\varphi\wedge\lambda_S^{-1}s\\\text{par définition}}}
\longrightarrow\lambda_{S^\vee}(\varphi(s))
\]

on voit que (76) est un isomorphisme de $Z$-ensembles. \deleted{\unsure{C}} Est-ce même un isomorphisme de $\unsure{\mathcal G|G}$-\deleted{torseurs} \inserted{ensembles} ? (En notant que du fait que $\mathcal G$ agit \deleted{\unsure{comme gr}} \inserted{sur $S!$} par $Z$-automorphismes, il agit aussi sur l'ensemble tordu $\Pi\wedge_Z(S!)$ par transport de structure, par

\[
\begin{aligned}
&g(\varphi\wedge s)=\varphi\wedge g_S(s)\qquad)\\
&\text{On aura}\quad g(\varphi\wedge s)=\varphi\wedge g_S(s)\longmapsto
\varphi(g_S(s))\overset?=g_{S^\vee}\varphi(s).
\end{aligned} \tag*{(77)}
\]

% archive sheet number: 216
\scanpage{218}
\deleted{\unsure{est}} \deleted{La relation signifie donc que pour} Ainsi

\[
\begin{aligned}
&\Pi\wedge_Z(S!)\xrightarrow{\sim} S^\vee\ \text{commute à un }g\in\unsure{\mathcal G|G}\ \text{donné}\\
&\text{si et seulement si pour tout }\varphi\in\Pi\ \text{-- ou}\\
&\text{simplement pour \uline{un} }\varphi_0\in\Pi\ \text{-- on a}\\
&\varphi(g_S(s))=g_{S^\vee}\varphi(s),\ \text{i.e. }\varphi\text{ commute à }g.
\end{aligned} \tag*{(78)}
\]

Donc, pour un sous-groupe donné $G\subset\mathcal G\subset\operatorname{Aut}(\Gamma)$ (on s'intéressera au cas où $Z\subset G$), l'isom. (76) est un isom de $G$-ens. ssi pour \uline{tout} $\varphi\in\Pi$ -- ou pour un $\varphi_0\in\Pi$, c'est kif-kif -- $\varphi$ commute à l'action de $\mathcal G$. \deleted{Le plus} Il y a un plus grand tel groupe $G$, qui \deleted{sont formé des $g\in\mathcal G$ qui} \inserted{est donné par}

\[
\mathcal G_0=\left\{g\in\operatorname{Aut}(S,S^\vee,R)\ \middle|\
\begin{aligned}
&\text{a) }g\text{ centralise }Z\ \text{i.e. }g\in\mathcal G\\
&\text{b) }\deleted{\unsure{\forall}}\exists\varphi\in\Pi,\ \varphi:S\xrightarrow{\sim}S^\vee\ \text{commute}\\
&\hspace{12mm}\text{à }g\ \text{i.e. }\varphi(g_Ss)=g_{S^\vee}(\varphi(s))
\end{aligned}\right\} \tag*{(79)}
\]

\margin{\uline{NB} Se garder de confondre cette condition de commutation au sens strict, avec la condition $g^\vee\varphi=\varphi g^{-1}$ (au lieu de $\varphi g$) !}

% slanted
\margin{Exemple un peu idiot, pas très instructif. Exemple important par contre : celui de $S,S^\vee$ déduits de $M^*,M^{\vee *}$ (p. 104' ff) \deleted{\unsure{alors}} Alors la \unsure{partie|trace} de $\mathcal G_0$ sur le groupe $\operatorname{Gl}_{\Lambda}^{\unsure{\prime|{}}}(M)$ $\mathcal G_0$ est formé des \unsure{il faut|u qui} $u\in\operatorname{Gl}_\Lambda(M)$ satisfaisant $\check u\varphi_0=\pm\varphi_0u$ i.e. $\underbrace{\check u\varphi_0u^{-1}}_{\varphi_0\det(u)}=\pm\varphi_0$ i.e. $\det u\in\pm1$. On a alors $(79\ \mathrm{bis})\quad\mathcal G_0=\operatorname{Gl}'_\Lambda(M)^{\pm1}$ du moins la \unsure{trace|prise|partie} du $\mathcal G_0$ \unsure{l'y|(79)} sur $\operatorname{Gl}'_\Lambda(M)$ \ldots}

\uline{Ex} Prenons $Z=1$ donc $\mathcal G=\operatorname{Aut}(\Gamma)$. Alors $\Pi$ se réduit à un élément, qui est n'importe quelle jauge $\varphi_0$, de $\Gamma=(S,S^\vee,R)$. Bien entendu, il n'est pas obligé que $\varphi_0$ commute à l'action de $\mathcal G$ -- p. ex. dans le cas de la p. 110 $S=S^\vee$, $R={}$ diagonale, $\varphi_0$ est donc une involution de $S$, tandis que $\mathcal G=\unsure{\mathfrak S|\mathscr S}_S$ ; ici $\mathcal G_0$ s'identifie au commutant de l'involution $\varphi_0$ dans $\mathfrak S_{E}$. \deleted{\unsure{Dans ce cas|Donc on a}}

Il faut encore voir comment se transforme par dualité une structure de multigraphe \inserted{(\deleted{\unsure{strict}} séparé)} à homothéties, \unsure{interprété|réinterprété} dans l'optique présente

% archive sheet number: 217
\scanpage{219}
\mspage{112}

comme une structure

\[
\bigl(S,S^\vee,R\subset S\times S^\vee,Z\subset\operatorname{Aut}(\underbrace{S\unsure{,|\times}S^\vee,R}_{\Gamma}),\Pi\subset\operatorname{Bij}(S,S^\vee)\bigr)\tag*{(81)}
\]

où, pour mémoire

\[
\left\{\begin{aligned}
&\text{a) } S,S^\vee\text{ deux ensembles, }R\subset S\times S^\vee\text{ une relation}\\
&\text{b) } Z\text{ un sous-groupe \uline{commutatif} de }\operatorname{Aut}(\Gamma)\\
&\text{c) } \Pi\text{ un ensemble de bijections }S\longrightarrow S^\vee,\text{ satisfaisant}\\
&\text{(1) condition de symétrie \unsure{ci-dessus|ci-dev.|cf.} \unsure{de|du} }Z\ \unsure{\ne|?}\ \unsure{\pi|\Pi|70}\\
&\text{(2) }\lambda_{S^\vee}\varphi=\varphi\lambda_S^{-1},\text{ i.e. }\lambda_{S^\vee}\varphi\lambda_S=\varphi\qquad\forall\lambda\in Z\\
&\deleted{\text{\unsure{i.e.} \unsure{cette condition|cette condit.}}}\\
&[\,\mathrm{NB}\ \text{les }\varphi\text{ satisfaisant à \unsure{ces|aux} conditions (1,2)}\\
&\qquad\text{forment un }\inserted{\text{ens }\Pi_Z\ \text{où }Z\text{ opère \uline{librement}}}\ \deleted{\text{ensemble}}\ \deleted{\unsure{G|g}\subset\operatorname{Aut}(\Gamma)}\\
&\qquad\text{par }T_g\varphi=g_{S^\vee}\varphi=\varphi\circ g_S^{-1}\,]\\
&\text{(3) }\Pi\text{ est un sous-}Z\text{-torseur de }\Pi_Z\\
&\deleted{\text{\unsure{cad} }\forall\varphi\in\Pi\ \text{(\unsure{$Z\ni$} pour un }\unsure{\varphi_0}\in\Pi,\text{ c'est kif-kif)}}
\end{aligned}\right.\tag*{(82)}
\]

% slanted
\margin{La condition (1), \unsure{assure|sous} que les données a) b) c), définissant un \inserted{ens. de} \deleted{parties} structures de \unsure{parties|graphes} \deleted{\unsure{strict}} strict séparés (moins \unsure{exiger} que $R$ soit séparante), ayant \deleted{un} \unsure{pr conséquence} \unsure{sous-entendu}, la condition (2) que $Z$ est un groupe d'homothéties pour cette structure de graphe.}

% slanted
\margin{\uline{NB} Il suffit de poser (2) pour \uline{un} $\varphi_0\in\Pi$, pour}

Si on considère l'objet

\[
\bigl(S'=S^\vee,S'^\vee=S,R'={}^sR,Z'={}^sZ,\Pi'=\Pi^{-1}=\{\varphi^{-1}\mid\varphi\in\Pi\}\bigr),\tag*{(83)}
\]

on voit qu'elle satisfait aux \deleted{\unsure{m}} conditions (82) que le système de départ (80). Notons que la bijection canonique

\[
\varphi\longmapsto\varphi^{-1}\qquad\Pi\xrightarrow{\sim}\Pi'\tag*{(84)}
\]

satisfait

\[
\bigl(\underbrace{T_\lambda\varphi}_{\lambda_{S^\vee}\varphi}\bigr)^{-1}
=\underbrace{T'_{\lambda'}(\varphi^{-1})}_{\deleted{\unsure{\lambda}}\ \varphi^{-1}\lambda_{S'}^{-1}}
\qquad\text{grand succès \unsure{!|?}},\tag*{(85)}
\]

% slanted
\margin{dans cette formule, j'identifie, \unsure{pr|je} \unsure{pos.|pose} $\lambda'=(\lambda_{S^\vee},\lambda_S)$ à $\lambda=(\lambda_S,\lambda_{S^\vee})$}

\margin{pour l'iso $\lambda\longmapsto\lambda'$, $Z\longrightarrow Z'$.}

i.e. (84) est \deleted{\unsure{équiv.}} un iso de $Z$-ensembles ! Si on considère d'autre part l'iso déduit de (76) en \deleted{\unsure{ten}} faisant le twist par \deleted{$\Pi^{\unsure{(-1)|\mathrm{inv}}}$} le torseur inverse de \deleted{$\Pi$}

\[
S!\simeq\Pi^{\mathrm{inv}}\wedge_Z S^\vee
\]

on en déduit, en passant à l'action opposée de $Z$ sur les deux mb

\[
S\simeq\Pi\wedge_Z(S^\vee!),\tag*{(86)}
\]

% archive sheet number: 218
